%{{{Header
\documentclass[a4paper,10pt]{article} \usepackage{anysize}
\marginsize{2cm}{2cm}{1cm}{1cm}
%\textwidth 6.0in \textheight = 664pt
\usepackage{xltxtra}
\usepackage{graphicx}
\usepackage{color}
\usepackage{xgreek}
\usepackage{fancyvrb}
\usepackage{minted}
\usepackage{listings}
\usepackage{enumitem}
\usepackage{framed}
\usepackage{relsize}
\usepackage{float}
\setmainfont[Mapping=TeX-text]{DejaVu Serif}
\newcommand{\tab}{\hspace*{3em}}

%\setmainfont{Arial}
%\setsansfont{FreeSans}
%\setmonofont{FreeMono}
\begin{document}
%\def\thesubsection {\alph{subsection}}
\renewcommand{\labelenumi}{\roman{enumi})}
\renewcommand{\labelenumii}{(\arabic{enumii})}

\begin{titlepage}
\begin{center}

\noindent\begin{minipage}[b]{.60\textwidth}
\begin{Large}
\emph{Εθνικό Μετσόβιο Πολυτεχνείο\\
Σχολή Ηλεκτρολόγων Μηχ. \& Μηχ. Υπολογιστών \\
Αλγόριθμοι και Πολυπλοκότητα\\
7\textsuperscript{ο} εξάμηνο, Ροή Λ\\
Ακαδημαϊκή περίοδος: 2010-2011\\}
\end{Large}
\end{minipage}%
\begin{minipage}[b]{.40\textwidth}
     \centering
     \includegraphics[scale=0.8]{title/ntua_logo}
\end{minipage}

\vspace{5cm}
\begin{Huge}
\underline{1\textsuperscript{η} Σειρά Γραπτών Ασκήσεων}
\end{Huge}
\vfill

\begin{flushright}
\Large{Γερακάρης Βασίλης}\\
\large{Α.Μ.: 03108092}\\
\end{flushright}
\vspace{1cm}

\large\today\\
\end{center}
\end{titlepage}


%}}}

%{{{Ταξινόμηση
\section{Ασυμπτωτικός συμβολισμός, Αναδρομικές
Σχέσεις} \setcounter{section}{1}
\begin{enumerate}

\item Ταξινόμηση
\begin{enumerate}
\item $\log {n^3} = 3 \log{n} = \Theta(logn)$
\item $\sqrt{n}*\log^{50}{n} $
\item $\frac{n}{\log{\log{n}} } $
\item $\log{n!} = \Theta (nlog{n})$
\item $n*\log^{10}{n} $
\item $n^{1.01} $
\item $5^{\log_{2}{n}} = n^{log_{2}{5}} = \Theta (n^{2.322}) $
\item $\sum_{k=1}^n{k^5} = {k^6}$
\item $\log^{\log{n}}{n} = n^{\log{\log{n}} } $
\item $2^{\log_2^4{n}} $
\item $\log^{\sqrt{n}}{n} $
\item $e^{\frac{n}{\ln{n}} } $
\item $n3^n $
\item $2^{2n} $
\item $\sqrt{n!} $
\end{enumerate}
%}}}

%{{{Τάξη μεγέθους
\item Τάξη Μεγέθους
\begin{enumerate}
\item $T(n)=5T(n/7)+n\log{n}$ ,\tab ($a=5 , b=7$) \\
$\Rightarrow n^{\log _7 {5}} = n^{0.827}$
$\Rightarrow n^{0.827} < n < n\log{n}$\\
$\Rightarrow T(n) \in \Theta(n\log{n})$ \tab M.Th (case \#3)\\


\item $T(n) = 4T(n/5)+n/\log^2{n}$ ,\tab ($a=4 , b=5$)\\
$\Rightarrow n^{\log_5 {4}} = n^{0.861}$
$\Rightarrow n^{0.861} < n/\log^2{n}$\\
$\Rightarrow T(n) \in \Theta(n/\log^2{n})$ \tab M.Th (case \#3)\\

\item $T(n) = T(n/3) +3T(n/7)+n$
\tab ($\frac{1}{3}+\frac{1}{7}+\frac{1}{7}+\frac{1}{7}< 1)$\\
$\Rightarrow T(n) \in \Theta(n)$ \tab M.Th (special case) \\

\item $T(n) = 6T(n/6)+n$\\
$\Rightarrow n^{\log _6 {6}} = n$\\
$\Rightarrow T(n) \in \Theta(n\log {n})$ \tab M.Th (sp. case \#2)\\

\item $T(n) = T(n/3) + T(2n/3) + n$\\
$\Rightarrow T(n) \in \Theta(n\log{n})$\\

\item $T(n) = 16T(n/4) + n^3\log^2{n}$ , \tab ($a=16, b=4)$\\
$\Rightarrow n^{\log _4{16}} = n^2$
$\Rightarrow n^2 < n^3\log^2{n}$\\
$\Rightarrow T(n) \in \Theta(n^3\log^2{n})$ \tab M.Th (case \#3)\\

\item $T(n) = T(\sqrt{n}) + \Theta(\log{\log{n}})$ \tab
(Θέτω $k=\log{n}, f(k)=T(n))$\\
$\Rightarrow
T(\sqrt{n})=f(\sqrt{k})=f(\log{\sqrt{n}})=f(\frac{\log{n}}{2})=f(\frac{k}{2}) $\\
$\Rightarrow f(k)=f(\frac{k}{2})+\Theta(\log{k})$\\
$\Rightarrow \Theta(n)$ \tab (M.Th (case \#1)\\

\item $T(n) = T(n-3) + \log{n}$ \tab ($\frac{n}{3} * \log{n})$ \\
$\Rightarrow T(n) \in \Theta(nlogn)$

\end{enumerate}
\end{enumerate}
\pagebreak
%}}}

%{{{Ταξινόμηση σε Πίνακα με Πολλά Ίδια Στοιχεία
\section{Ταξινόμηση σε Πίνακα με Πολλά Ίδια Στοιχεία}

Δημιουργούμε ένα 2-διάστατο πίνακα $B[\log ^d {n},2]$ όπου η 1η γραμμή
θα περιέχει τα διαφορετικά στοιχεία ταξινομημένα και η 2η το πλήθος
εμφάνισης του κάθε στοιχείου.\\
Αρχίζουμε θέτονταςτο $B[1,1] = A[1]$, $B[1,2] = 1$, i=2.\\
Για κάθε $A[i]$ στον πίνακα, εφαρμόζουμε Binary Search στην πρώτη γραμμή του
πίνακα Β. \\
(Με κόστος $\Theta (n log{log^d{n}})$ )
\begin{itemize}
\item Αν το στοιχείο δεν υπάρχει στον πίνακα B (H binary search επιστρέψει 0)\\
τότε εφαρμόζουμε Insertion sort του στοιχείου στον πίνακα Β, τοποθετώντας το
στη σωστή θέση
(αντιμεταθέτωντας και τη 2η γραμμή όταν χρειάζεται)\\
Με κόστος: $(\log ^d{n})^2$
\item Αν το στοιχείο υπάρχει στον πίνακα (έστω στη θέση k), αυξάνουμε το
$B[k,2]$ κατά 1.
\end{itemize}
Όταν τελειώσουμε, αναπτύσσουμε τα δεδομένα του (ταξινομημένου) πίνακα Β (το στοιχείο $B[i,1]$
αναπτύσσεται $B[i,2]$ φορές) σε γραμμικό χρόνο.\\

Συνολικό κόστος: $(n log{log^d{n}}) + (\log ^d{n})^2 + n =  O(n log{log^d{n}}) $

%}}}

%{{{Δυαδική Αναζήτηση
\section{Δυαδική Αναζήτηση}
\begin{enumerate}
\item Ξεκινάμε (για $i=1$) συγκρίνοντας το $A[i]$ με το $x$. Όσο το $x$ είναι
μεγαλύτερος από αυτό, κρατάμε την τιμή του $(i+1)$ σε μια άλλη μεταβλητή $k$ και
διπλασιάζουμε το i.
Μόλις φτάσουμε σε μεγαλύτερο αριθμό από τον $x$, εφαρμόζω ένα έλεγχο.
Αν το στοιχείο $A[i]$ είναι διάφορο του $\infty$, τότε εφαρμόζω δυαδική αναζήτηση
για το στοιχείο $x$ στο τμήμα $Α[k..i]$ του πίνακα και έχω το αποτέλεσμα για
πιθανή ύπαρξη \& θέση του στοιχείου. \\
Αν το στοιχείο $A[i]$ είναι ίσο με $\infty$, τότε ελέγχω αν $k=i$ (τερματική
συνθήκη) και αν όχι ελέγχω το $A[(k+i) div 2]$.
\begin{itemize}
\item Αν είναι $\infty$ τότε θέτω $i= (k+i)div 2$ και επαναλαβάνω.
\item Αν είναι αριθμός $m$ τον συγκρίνω με το $x$.
\begin{itemize}
\item Αν $x>m$ τότε θέτω $k= (k+i) div 2$ και επαναλαμβάνω.
\item Αν $x<m$ τότε εφαρμόζω δυαδική αναζήτηση στο τμήμα $A[k..i]$ του
πίνακα.
\item Αν $x=m$ τότε είναι το ζητούμενο στοιχείο.
\end{itemize}
\end{itemize}

Η τελική πολυπλοκότητα είναι $O(\log{n})$.\\

\item Εργαζόμαστε με τα k πρώτα στοιχεία από κάθε πίνακα (συμπληρώνοντας τις
κενές θέσεις με + $\infty$ , αν ο πίνακας δεν έχει τόσα στοιχεία). Συγκρίνουμε τα
$A[k/2]$ με $B[k/2+1]$ και $A[k/2+1]$ με $B[k/2]$.
\begin{itemize}
\item Αν $A[k/2]<B[k/2+1]$ και $B[k/2]<A[k/2+1]$ τοτε το k-οστό στοιχείο είναι το
$max\{A[k/2],B[k/2]\}$.
\item Αν $A[k/2]<B[k/2+1]$ και $B[k/2]>A[k/2+1]$ τότε επαναλαμβάνουμε αναδρομικά
τη διαδικασία στους υποπίνακες $A[k/2+1]..A[k]$ και $B[1]..B[k/2]$ για $k = \frac
{k}{2}$.
\end{itemize}

Τελική πολυπλοκότητα: $O(log{k})$
\end{enumerate}
%}}}

%{{{Συλλογή Comics
\section{Συλλογή Comics}
Ελέγχουμε αρχικά το 1\textsuperscript{o} bit για όλα τα τεύχη, χωρίζοντάς τα
σε 2 σύνολα με πλήθος στοιχείων $\frac{n}{2}$ και $\frac{n}{2}-1$. Το τέυχος
που λείπει είναι στο σύνολο με τα λιγότερα στοιχεία, οπότε αποθηκεύουμε την
τιμή αυτού του bit. \\
Εφαρμόζουμε αναδρομικά τη μέθοδο για το επόμενο bit, κάθε φορά στο σύνολο με τα λιγότερα στοιχεία,
κάνοντας τις μισές ερωτήσεις από το προηγούμενο βήμα σε κάθε αναδρομή. \\
Τελικά, θα καταλήξουμε στον κωδικό του τεύχους που λείπει από τη συλλογή, έχοντας
ρωτήσει $n+\frac{n}{2}+\frac{n}{4}+..+1 \approx 2n$ φορές.\\
Πρακτικά, είναι: $(n-1)+(\frac{n}{2}-1)+(\frac{n}{4} -1)+..=
\displaystyle\sum\limits_{i=0}^{\log{n}} (\frac{n}{2^i} -1) = 2n - log{n} -2$


\pagebreak
%}}}

%{{{Πολυκατοικίες χωρίς Θέα
\section{Πολυκατοικίες χωρίς Θέα}
Αρχίζουμε αρχικοποιώντας $B[1] = 0$, μιάς και ο 1\textsuperscript{ος} δεν έχει
κτήριο δυτικά του.\\
Για κάθε κτήριο i ανατολικά του, θέτουμε $k= i-1$ και συγκρίνουμε το ύψος του $A[i]$ με του
$A[k]$.
\begin{itemize}
\item Αν $A[i]<A[k]$, τότε $B[i] = k$ και συνέχιζουμε στο επόμενο δεξιά
κτήριο ($i= i+1$) .
\item Αν $A[i] \ge A[k]$,
\begin{itemize}
\item Αν $B[k] = 0$ τότε $B[i] = 0$ και συνέχιζουμε στο επόμενο δεξιά
κτήριο.
\item Αλλιώς θέτουμε $k = B[k]$ και επαναλαμβάνουμε τον έλεγχο.
\end{itemize}
\end{itemize}

%}}}

\end{document}
